% Options for packages loaded elsewhere
\PassOptionsToPackage{unicode}{hyperref}
\PassOptionsToPackage{hyphens}{url}
\documentclass[
  11pt,
  a4paper,
]{article}
\usepackage{xcolor}
\usepackage[margin=18mm]{geometry}
\usepackage{amsmath,amssymb}
\setcounter{secnumdepth}{-\maxdimen} % remove section numbering
\usepackage{iftex}
\ifPDFTeX
  \usepackage[T1]{fontenc}
  \usepackage[utf8]{inputenc}
  \usepackage{textcomp} % provide euro and other symbols
\else % if luatex or xetex
  \usepackage{unicode-math} % this also loads fontspec
  \defaultfontfeatures{Scale=MatchLowercase}
  \defaultfontfeatures[\rmfamily]{Ligatures=TeX,Scale=1}
\fi
\ifPDFTeX\else
  % xetex/luatex font selection
    \setmainfont[]{Noto Serif}
    \setsansfont[]{Noto Sans}
    \setmonofont[]{Noto Sans Mono}
  \setmathfont[]{Noto Sans Math}
\fi
% Use upquote if available, for straight quotes in verbatim environments
\IfFileExists{upquote.sty}{\usepackage{upquote}}{}
\IfFileExists{microtype.sty}{% use microtype if available
  \usepackage[]{microtype}
  \UseMicrotypeSet[protrusion]{basicmath} % disable protrusion for tt fonts
}{}
\makeatletter
\@ifundefined{KOMAClassName}{% if non-KOMA class
  \IfFileExists{parskip.sty}{%
    \usepackage{parskip}
  }{% else
    \setlength{\parindent}{0pt}
    \setlength{\parskip}{6pt plus 2pt minus 1pt}}
}{% if KOMA class
  \KOMAoptions{parskip=half}}
\makeatother
\ifLuaTeX
\usepackage[bidi=basic]{babel}
\else
\usepackage[bidi=default]{babel}
\fi
\babelprovide[main,import]{russian}
\ifPDFTeX
\else
\babelfont{rm}[]{Noto Serif}
\fi
% get rid of language-specific shorthands (see #6817):
\let\LanguageShortHands\languageshorthands
\def\languageshorthands#1{}
\setlength{\emergencystretch}{3em} % prevent overfull lines
\providecommand{\tightlist}{%
  \setlength{\itemsep}{0pt}\setlength{\parskip}{0pt}}
\usepackage{bookmark}
\IfFileExists{xurl.sty}{\usepackage{xurl}}{} % add URL line breaks if available
\urlstyle{same}
\hypersetup{
  pdftitle={Билеты по дискретным случайным процессам},
  pdflang={ru-RU},
  hidelinks,
  pdfcreator={LaTeX via pandoc}}

\title{Билеты по дискретным случайным процессам}
\author{}
\date{}

\begin{document}
\maketitle

\section{Билет 18. Гауссовы случайные величины, векторы и
последовательности}\label{ux431ux438ux43bux435ux442-18.-ux433ux430ux443ux441ux441ux43eux432ux44b-ux441ux43bux443ux447ux430ux439ux43dux44bux435-ux432ux435ux43bux438ux447ux438ux43dux44b-ux432ux435ux43aux442ux43eux440ux44b-ux438-ux43fux43eux441ux43bux435ux434ux43eux432ux430ux442ux435ux43bux44cux43dux43eux441ux442ux438}

Источники: Ширяев 1, §13; Ширяев 2, гл. VI.

\subsection{1. Короткий
ответ}\label{ux43aux43eux440ux43eux442ux43aux438ux439-ux43eux442ux432ux435ux442}

Нормальная случайная величина с параметрами \(m\in\mathbb R\) и
\(\sigma^2\ge0\) имеет характеристическую функцию

\[
\varphi_X(t)
=
\mathbb E e^{itX}
=
\exp\left\{itm-\frac12\sigma^2t^2\right\}.
\]

Если \(\sigma^2>0\), то ее плотность:

\[
p(x)=\frac1{\sqrt{2\pi\sigma^2}}
\exp\left\{-\frac{(x-m)^2}{2\sigma^2}\right\}.
\]

Случайный вектор

\[
X=(X_1,\ldots,X_n)^T
\]

называется гауссовским, если любая его линейная комбинация

\[
\sum_{j=1}^n u_jX_j
\]

имеет нормальное, возможно вырожденное, распределение. Закон
гауссовского вектора полностью задается средним вектором

\[
m=\mathbb EX
\]

и ковариационной матрицей

\[
\Sigma=\operatorname{Cov}(X).
\]

Его характеристическая функция:

\[
\varphi_X(u)
=
\exp\left\{i\langle u,m\rangle-\frac12u^T\Sigma u\right\}.
\]

Случайная последовательность называется гауссовской, если любой ее
конечный вектор

\[
(X_{n_1},\ldots,X_{n_k})
\]

является гауссовским.

Главные свойства: линейные преобразования гауссовских векторов гауссовы;
для совместно гауссовских величин некоррелированность эквивалентна
независимости; для гауссовской последовательности широкая стационарность
влечет узкую.

\subsection{2. Полный
ответ}\label{ux43fux43eux43bux43dux44bux439-ux43eux442ux432ux435ux442}

\subsubsection{2.1. Введение в гауссовские распределения и
процессы}\label{ux432ux432ux435ux434ux435ux43dux438ux435-ux432-ux433ux430ux443ux441ux441ux43eux432ux441ux43aux438ux435-ux440ux430ux441ux43fux440ux435ux434ux435ux43bux435ux43dux438ux44f-ux438-ux43fux440ux43eux446ux435ux441ux441ux44b}

Гауссовы случайные величины и процессы важны потому, что они являются
основным классом распределений второго порядка. Для произвольного
процесса среднее и ковариация обычно не задают весь закон. Для
гауссовского процесса задают: все конечномерные распределения полностью
определяются средними и ковариациями.

\subsubsection{2.2. Одномерное нормальное
распределение}\label{ux43eux434ux43dux43eux43cux435ux440ux43dux43eux435-ux43dux43eux440ux43cux430ux43bux44cux43dux43eux435-ux440ux430ux441ux43fux440ux435ux434ux435ux43bux435ux43dux438ux435}

Начнем с одномерного случая. Случайная величина \(X\) имеет нормальное
распределение с параметрами \(m\) и \(\sigma^2>0\), если ее плотность
равна

\[
p(x)
=
\frac1{\sqrt{2\pi\sigma^2}}
\exp\left\{-\frac{(x-m)^2}{2\sigma^2}\right\}.
\]

Пишут:

\[
X\sim N(m,\sigma^2).
\]

Здесь

\[
\mathbb EX=m,
\qquad
\operatorname{Var}X=\sigma^2.
\]

Стандартная нормальная величина имеет распределение

\[
N(0,1).
\]

Если \(Z\sim N(0,1)\), то

\[
X=m+\sigma Z
\]

имеет распределение \(N(m,\sigma^2)\).

\subsubsection{2.3. Вырожденный случай и характеристическая
функция}\label{ux432ux44bux440ux43eux436ux434ux435ux43dux43dux44bux439-ux441ux43bux443ux447ux430ux439-ux438-ux445ux430ux440ux430ux43aux442ux435ux440ux438ux441ux442ux438ux447ux435ux441ux43aux430ux44f-ux444ux443ux43dux43aux446ux438ux44f}

Разрешают и вырожденный случай \(\sigma^2=0\). Тогда

\[
X=m
\]

почти наверное. Это тоже удобно считать нормальным распределением: оно
не имеет обычной плотности относительно меры Лебега, но его
характеристическая функция получается из той же формулы.

Характеристическая функция нормальной величины:

\[
\varphi_X(t)
=
\mathbb E e^{itX}
=
\exp\left\{itm-\frac12\sigma^2t^2\right\}.
\]

Эта формула важнее плотности, потому что она сразу обобщается на
гауссовские векторы.

\subsubsection{2.4. Определение гауссовского
вектора}\label{ux43eux43fux440ux435ux434ux435ux43bux435ux43dux438ux435-ux433ux430ux443ux441ux441ux43eux432ux441ux43aux43eux433ux43e-ux432ux435ux43aux442ux43eux440ux430}

Теперь случайный вектор. Вектор

\[
X=(X_1,\ldots,X_n)^T
\]

называется гауссовским, если для любого вектора коэффициентов

\[
u=(u_1,\ldots,u_n)^T\in\mathbb R^n
\]

линейная комбинация

\[
\langle u,X\rangle
=
\sum_{j=1}^n u_jX_j
\]

имеет нормальное распределение, возможно вырожденное.

Важно: недостаточно, чтобы каждая координата \(X_j\) отдельно была
нормально распределена. Нужно, чтобы нормальной была любая линейная
комбинация координат. Именно это условие контролирует совместное
распределение.

\subsubsection{2.5. Параметры и распределение гауссовского
вектора}\label{ux43fux430ux440ux430ux43cux435ux442ux440ux44b-ux438-ux440ux430ux441ux43fux440ux435ux434ux435ux43bux435ux43dux438ux435-ux433ux430ux443ux441ux441ux43eux432ux441ux43aux43eux433ux43e-ux432ux435ux43aux442ux43eux440ux430}

Для гауссовского вектора определены:

\[
m=\mathbb EX
=
(\mathbb EX_1,\ldots,\mathbb EX_n)^T
\]

и ковариационная матрица

\[
\Sigma=(\Sigma_{ij}),
\qquad
\Sigma_{ij}
=
\operatorname{Cov}(X_i,X_j).
\]

Матрица \(\Sigma\) симметрична и неотрицательно определена:

\[
u^T\Sigma u
=
\operatorname{Var}\langle u,X\rangle
\ge0.
\]

Характеристическая функция гауссовского вектора имеет вид

\[
\varphi_X(u)
=
\mathbb E e^{i\langle u,X\rangle}
=
\exp\left\{
i\langle u,m\rangle
-\frac12u^T\Sigma u
\right\}.
\]

Эта формула показывает, что закон гауссовского вектора полностью
задается парой

\[
(m,\Sigma).
\]

Обратно, если \(m\in\mathbb R^n\), а \(\Sigma\) - симметричная
неотрицательно определенная матрица, то существует гауссовский вектор с
такими средним и ковариацией. Если \(\Sigma\) положительно определена,
то распределение имеет плотность:

\[
p(x)
=
\frac1{(2\pi)^{n/2}(\det\Sigma)^{1/2}}
\exp\left\{
-\frac12(x-m)^T\Sigma^{-1}(x-m)
\right\}.
\]

Если \(\Sigma\) вырождена, то плотности относительно \(n\)-мерной меры
Лебега нет: распределение сосредоточено на некотором аффинном
подпространстве. Но вектор все равно называется гауссовским, потому что
все линейные комбинации нормальны.

\subsubsection{2.6. Линейные преобразования гауссовских
векторов}\label{ux43bux438ux43dux435ux439ux43dux44bux435-ux43fux440ux435ux43eux431ux440ux430ux437ux43eux432ux430ux43dux438ux44f-ux433ux430ux443ux441ux441ux43eux432ux441ux43aux438ux445-ux432ux435ux43aux442ux43eux440ux43eux432}

Линейные преобразования гауссовских векторов снова гауссовы. Если

\[
Y=AX+b,
\]

где \(A\) - матрица, \(b\) - вектор, а \(X\) гауссовский, то \(Y\)
гауссовский. Его параметры:

\[
\mathbb EY=Am+b,
\qquad
\operatorname{Cov}(Y)=A\Sigma A^T.
\]

Доказательство простое: любая линейная комбинация координат \(Y\)
является линейной комбинацией координат \(X\) плюс константа, а значит
нормальна.

\subsubsection{2.7. Свойства независимости и
некоррелированности}\label{ux441ux432ux43eux439ux441ux442ux432ux430-ux43dux435ux437ux430ux432ux438ux441ux438ux43cux43eux441ux442ux438-ux438-ux43dux435ux43aux43eux440ux440ux435ux43bux438ux440ux43eux432ux430ux43dux43dux43eux441ux442ux438}

Главное специальное свойство гауссовских векторов: для совместно
гауссовских случайных величин некоррелированность равносильна
независимости. В общем случае из

\[
\operatorname{Cov}(X,Y)=0
\]

независимость не следует, как обсуждалось в Билете 14. Но если \((X,Y)\)
- гауссовский вектор, то из нулевой ковариации следует независимость.

Более общая формулировка: если гауссовский вектор разбит на две группы
координат

\[
X=(X^{(1)},X^{(2)}),
\]

и все ковариации между первой и второй группой равны нулю, то эти два
случайных вектора независимы. В терминах матрицы ковариаций это означает
блочно-диагональную матрицу:

\[
\Sigma=
\begin{pmatrix}
\Sigma_1 & 0\\
0 & \Sigma_2
\end{pmatrix}.
\]

Тогда характеристическая функция факторизуется:

\[
\varphi_X(u_1,u_2)
=
\varphi_{X^{(1)}}(u_1)\varphi_{X^{(2)}}(u_2),
\]

а факторизация характеристической функции означает независимость.

\subsubsection{2.8. Гауссовские
последовательности}\label{ux433ux430ux443ux441ux441ux43eux432ux441ux43aux438ux435-ux43fux43eux441ux43bux435ux434ux43eux432ux430ux442ux435ux43bux44cux43dux43eux441ux442ux438}

Теперь гауссовские последовательности. Случайная последовательность

\[
\{X_n:n\in T\}
\]

называется гауссовской, если для любого конечного набора индексов

\[
n_1,\ldots,n_k
\]

случайный вектор

\[
(X_{n_1},\ldots,X_{n_k})
\]

является гауссовским. Эквивалентно: любая конечная линейная комбинация

\[
\sum_{j=1}^k u_jX_{n_j}
\]

имеет нормальное распределение.

Гауссовская последовательность задается средней функцией

\[
m_n=\mathbb EX_n
\]

и ковариационной функцией

\[
K(n,m)=\operatorname{Cov}(X_n,X_m).
\]

Для любого конечного набора индексов \(n_1,\ldots,n_k\) средний вектор
равен

\[
(m_{n_1},\ldots,m_{n_k}),
\]

а ковариационная матрица равна

\[
(K(n_i,n_j))_{i,j=1}^k.
\]

Чтобы такая последовательность существовала, функция \(K\) должна быть
симметричной и неотрицательно определенной:

\[
\sum_{i,j=1}^k a_i a_jK(n_i,n_j)\ge0
\]

для любых \(k\), индексов \(n_i\) и коэффициентов \(a_i\). Тогда
согласованные гауссовские конечномерные распределения задают процесс по
теореме Колмогорова из Билета 15.

\subsubsection{2.9. Стационарность гауссовских
последовательностей}\label{ux441ux442ux430ux446ux438ux43eux43dux430ux440ux43dux43eux441ux442ux44c-ux433ux430ux443ux441ux441ux43eux432ux441ux43aux438ux445-ux43fux43eux441ux43bux435ux434ux43eux432ux430ux442ux435ux43bux44cux43dux43eux441ux442ux435ux439}

Для стационарности гауссовских последовательностей есть важное
упрощение. В Билете 17 различались узкая и широкая стационарность. Для
произвольных процессов широкая стационарность не влечет узкую. Но для
гауссовской последовательности влечет. Если

\[
\mathbb EX_n=m
\]

и

\[
\operatorname{Cov}(X_n,X_{n+k})=R(k),
\]

то все конечномерные распределения не меняются при сдвиге, потому что у
них не меняются средние векторы и ковариационные матрицы. Значит
гауссовская последовательность стационарна в узком смысле.

\subsection{3. Основные
определения}\label{ux43eux441ux43dux43eux432ux43dux44bux435-ux43eux43fux440ux435ux434ux435ux43bux435ux43dux438ux44f}

\subsubsection{Нормальная случайная
величина}\label{ux43dux43eux440ux43cux430ux43bux44cux43dux430ux44f-ux441ux43bux443ux447ux430ux439ux43dux430ux44f-ux432ux435ux43bux438ux447ux438ux43dux430}

Случайная величина \(X\) имеет распределение \(N(m,\sigma^2)\), если при
\(\sigma^2>0\) ее плотность равна

\[
p(x)=\frac1{\sqrt{2\pi\sigma^2}}
\exp\left\{-\frac{(x-m)^2}{2\sigma^2}\right\}.
\]

При \(\sigma^2=0\) понимают вырожденное распределение \(X=m\) почти
наверное.

\subsubsection{Характеристическая функция нормальной
величины}\label{ux445ux430ux440ux430ux43aux442ux435ux440ux438ux441ux442ux438ux447ux435ux441ux43aux430ux44f-ux444ux443ux43dux43aux446ux438ux44f-ux43dux43eux440ux43cux430ux43bux44cux43dux43eux439-ux432ux435ux43bux438ux447ux438ux43dux44b}

\[
\varphi_X(t)=\exp\left\{itm-\frac12\sigma^2t^2\right\}.
\]

\subsubsection{Гауссовский
вектор}\label{ux433ux430ux443ux441ux441ux43eux432ux441ux43aux438ux439-ux432ux435ux43aux442ux43eux440}

Вектор \(X=(X_1,\ldots,X_n)^T\) называется гауссовским, если для любого
\(u\in\mathbb R^n\) величина \(\langle u,X\rangle\) нормально
распределена, возможно вырожденно.

\subsubsection{Средний
вектор}\label{ux441ux440ux435ux434ux43dux438ux439-ux432ux435ux43aux442ux43eux440}

\[
m=\mathbb EX.
\]

\subsubsection{Ковариационная
матрица}\label{ux43aux43eux432ux430ux440ux438ux430ux446ux438ux43eux43dux43dux430ux44f-ux43cux430ux442ux440ux438ux446ux430}

\[
\Sigma_{ij}=\operatorname{Cov}(X_i,X_j).
\]

\subsubsection{Характеристическая функция гауссовского
вектора}\label{ux445ux430ux440ux430ux43aux442ux435ux440ux438ux441ux442ux438ux447ux435ux441ux43aux430ux44f-ux444ux443ux43dux43aux446ux438ux44f-ux433ux430ux443ux441ux441ux43eux432ux441ux43aux43eux433ux43e-ux432ux435ux43aux442ux43eux440ux430}

\[
\varphi_X(u)
=
\exp\left\{
i\langle u,m\rangle-\frac12u^T\Sigma u
\right\}.
\]

\subsubsection{Гауссовская
последовательность}\label{ux433ux430ux443ux441ux441ux43eux432ux441ux43aux430ux44f-ux43fux43eux441ux43bux435ux434ux43eux432ux430ux442ux435ux43bux44cux43dux43eux441ux442ux44c}

Последовательность \((X_n)\) называется гауссовской, если каждый
конечный вектор

\[
(X_{n_1},\ldots,X_{n_k})
\]

гауссовский.

\subsubsection{Белый гауссовский
шум}\label{ux431ux435ux43bux44bux439-ux433ux430ux443ux441ux441ux43eux432ux441ux43aux438ux439-ux448ux443ux43c}

Последовательность независимых одинаково распределенных гауссовских
величин с нулевым средним:

\[
X_n\sim N(0,\sigma^2),
\qquad
\operatorname{Cov}(X_n,X_m)=0\quad(n\ne m).
\]

\subsection{4. Основные теоремы и
свойства}\label{ux43eux441ux43dux43eux432ux43dux44bux435-ux442ux435ux43eux440ux435ux43cux44b-ux438-ux441ux432ux43eux439ux441ux442ux432ux430}

\begin{itemize}
\tightlist
\item
  Нормальный закон задается средним и дисперсией:
\end{itemize}

\[
X\sim N(m,\sigma^2).
\]

\begin{itemize}
\tightlist
\item
  Гауссовский вектор задается средним вектором и ковариационной
  матрицей:
\end{itemize}

\[
\varphi_X(u)
=
\exp\left\{i\langle u,m\rangle-\frac12u^T\Sigma u\right\}.
\]

\begin{itemize}
\tightlist
\item
  Матрица \(\Sigma\) должна быть симметричной и неотрицательно
  определенной.
\item
  Если \(\Sigma>0\), то гауссовский вектор имеет плотность:
\end{itemize}

\[
p(x)
=
\frac1{(2\pi)^{n/2}(\det\Sigma)^{1/2}}
\exp\left\{-\frac12(x-m)^T\Sigma^{-1}(x-m)\right\}.
\]

\begin{itemize}
\tightlist
\item
  Линейное преобразование гауссовского вектора гауссово:
\end{itemize}

\[
Y=AX+b
\quad\Rightarrow\quad
Y\ \text{гауссовский}.
\]

\begin{itemize}
\tightlist
\item
  Для совместно гауссовских величин некоррелированность равносильна
  независимости.
\item
  Нормальность отдельных координат не гарантирует совместную
  гауссовость.
\item
  Гауссовская последовательность задается средней функцией \(m_n\) и
  ковариационной функцией \(K(n,m)\).
\item
  Для гауссовской последовательности широкая стационарность влечет
  узкую.
\end{itemize}

\subsection{5. Примеры}\label{ux43fux440ux438ux43cux435ux440ux44b}

\subsubsection{Стандартная нормальная
величина}\label{ux441ux442ux430ux43dux434ux430ux440ux442ux43dux430ux44f-ux43dux43eux440ux43cux430ux43bux44cux43dux430ux44f-ux432ux435ux43bux438ux447ux438ux43dux430}

Если

\[
Z\sim N(0,1),
\]

то

\[
\varphi_Z(t)=e^{-t^2/2}.
\]

Если

\[
X=m+\sigma Z,
\]

то

\[
X\sim N(m,\sigma^2).
\]

\subsubsection{Двумерный гауссовский
вектор}\label{ux434ux432ux443ux43cux435ux440ux43dux44bux439-ux433ux430ux443ux441ux441ux43eux432ux441ux43aux438ux439-ux432ux435ux43aux442ux43eux440}

Пусть

\[
X=(X_1,X_2)^T
\]

имеет средний вектор

\[
m=(m_1,m_2)^T
\]

и ковариационную матрицу

\[
\Sigma=
\begin{pmatrix}
\sigma_1^2 & \rho\sigma_1\sigma_2\\
\rho\sigma_1\sigma_2 & \sigma_2^2
\end{pmatrix},
\qquad
|\rho|\le1.
\]

Если вектор гауссовский, то при \(\rho=0\) координаты независимы. Это
специальное гауссовское свойство.

\subsubsection{Нормальные координаты без совместной
гауссовости}\label{ux43dux43eux440ux43cux430ux43bux44cux43dux44bux435-ux43aux43eux43eux440ux434ux438ux43dux430ux442ux44b-ux431ux435ux437-ux441ux43eux432ux43cux435ux441ux442ux43dux43eux439-ux433ux430ux443ux441ux441ux43eux432ux43eux441ux442ux438}

Пусть \(Z\sim N(0,1)\), а \(\varepsilon\) независимо принимает значения
\(1\) и \(-1\) с вероятностями \(1/2\). Положим

\[
X=Z,
\qquad
Y=\varepsilon Z.
\]

Тогда \(X\) и \(Y\) по отдельности имеют распределение \(N(0,1)\). Но
вектор \((X,Y)\) не является гауссовским: например, сумма

\[
X+Y=(1+\varepsilon)Z
\]

с вероятностью \(1/2\) равна \(0\), а с вероятностью \(1/2\) равна
\(2Z\), поэтому не имеет нормального распределения. Этот пример
показывает, почему в определении нужна нормальность всех линейных
комбинаций, а не только координат.

\subsubsection{Белый гауссовский
шум}\label{ux431ux435ux43bux44bux439-ux433ux430ux443ux441ux441ux43eux432ux441ux43aux438ux439-ux448ux443ux43c-1}

Пусть

\[
X_n
\]

независимы и

\[
X_n\sim N(0,\sigma^2).
\]

Тогда \((X_n)\) - гауссовская последовательность,

\[
\mathbb EX_n=0,
\]

и

\[
R(k)=\operatorname{Cov}(X_0,X_k)
=
\begin{cases}
\sigma^2, & k=0,\\
0, & k\ne0.
\end{cases}
\]

Это базовый пример стационарной гауссовской последовательности.

\subsubsection{Гауссовский
AR(1)}\label{ux433ux430ux443ux441ux441ux43eux432ux441ux43aux438ux439-ar1}

Пусть

\[
X_n=aX_{n-1}+\xi_n,
\qquad |a|<1,
\]

где \(\xi_n\) - белый гауссовский шум с дисперсией \(\sigma_\xi^2\). В
стационарном режиме

\[
\mathbb EX_n=0,
\qquad
R(k)=\frac{\sigma_\xi^2}{1-a^2}a^{|k|}.
\]

Такой процесс гауссовский, потому что строится линейными операциями из
гауссовских величин.

\subsection{6. Схемы доказательств /
обоснований}\label{ux441ux445ux435ux43cux44b-ux434ux43eux43aux430ux437ux430ux442ux435ux43bux44cux441ux442ux432-ux43eux431ux43eux441ux43dux43eux432ux430ux43dux438ux439}

\subsubsection{\texorpdfstring{Почему закон гауссовского вектора
задается \(m\) и
\(\Sigma\)}{Почему закон гауссовского вектора задается m и \textbackslash Sigma}}\label{ux43fux43eux447ux435ux43cux443-ux437ux430ux43aux43eux43d-ux433ux430ux443ux441ux441ux43eux432ux441ux43aux43eux433ux43e-ux432ux435ux43aux442ux43eux440ux430-ux437ux430ux434ux430ux435ux442ux441ux44f-m-ux438-sigma}

Для любого \(u\) величина

\[
\langle u,X\rangle
\]

нормальна. Ее среднее:

\[
\mathbb E\langle u,X\rangle=\langle u,m\rangle.
\]

Ее дисперсия:

\[
\operatorname{Var}\langle u,X\rangle=u^T\Sigma u.
\]

Значит

\[
\mathbb E e^{i\langle u,X\rangle}
=
\exp\left\{i\langle u,m\rangle-\frac12u^T\Sigma u\right\}.
\]

Характеристическая функция однозначно задает распределение, поэтому
закон определяется \(m\) и \(\Sigma\).

\subsubsection{Почему линейное преобразование гауссовского вектора
гауссово}\label{ux43fux43eux447ux435ux43cux443-ux43bux438ux43dux435ux439ux43dux43eux435-ux43fux440ux435ux43eux431ux440ux430ux437ux43eux432ux430ux43dux438ux435-ux433ux430ux443ux441ux441ux43eux432ux441ux43aux43eux433ux43e-ux432ux435ux43aux442ux43eux440ux430-ux433ux430ux443ux441ux441ux43eux432ux43e}

Пусть

\[
Y=AX+b.
\]

Для любого \(v\):

\[
\langle v,Y\rangle
=
\langle v,AX+b\rangle
=
\langle A^Tv,X\rangle+\langle v,b\rangle.
\]

Это линейная комбинация координат \(X\) плюс константа, значит
нормальная величина. Следовательно, \(Y\) гауссовский.

\subsubsection{Почему некоррелированность дает независимость для
Гаусса}\label{ux43fux43eux447ux435ux43cux443-ux43dux435ux43aux43eux440ux440ux435ux43bux438ux440ux43eux432ux430ux43dux43dux43eux441ux442ux44c-ux434ux430ux435ux442-ux43dux435ux437ux430ux432ux438ux441ux438ux43cux43eux441ux442ux44c-ux434ux43bux44f-ux433ux430ux443ux441ux441ux430}

Пусть \(X=(X^{(1)},X^{(2)})\) - гауссовский вектор, и ковариации между
группами равны нулю. Тогда

\[
\Sigma=
\begin{pmatrix}
\Sigma_1 & 0\\
0 & \Sigma_2
\end{pmatrix}.
\]

Характеристическая функция:

\[
\varphi(u_1,u_2)
=
\exp\left\{
i\langle u_1,m_1\rangle+i\langle u_2,m_2\rangle
-\frac12u_1^T\Sigma_1u_1
-\frac12u_2^T\Sigma_2u_2
\right\}.
\]

Она раскладывается в произведение:

\[
\varphi(u_1,u_2)
=
\varphi_1(u_1)\varphi_2(u_2).
\]

Факторизация характеристической функции означает независимость.

\subsubsection{Почему широкая стационарность гауссовской
последовательности влечет
узкую}\label{ux43fux43eux447ux435ux43cux443-ux448ux438ux440ux43eux43aux430ux44f-ux441ux442ux430ux446ux438ux43eux43dux430ux440ux43dux43eux441ux442ux44c-ux433ux430ux443ux441ux441ux43eux432ux441ux43aux43eux439-ux43fux43eux441ux43bux435ux434ux43eux432ux430ux442ux435ux43bux44cux43dux43eux441ux442ux438-ux432ux43bux435ux447ux435ux442-ux443ux437ux43aux443ux44e}

Для любого конечного набора индексов \(n_1,\ldots,n_k\) распределение
вектора

\[
(X_{n_1},\ldots,X_{n_k})
\]

гауссовское и задается средним вектором и ковариационной матрицей. Если

\[
\mathbb EX_n=m,
\qquad
\operatorname{Cov}(X_n,X_{n+k})=R(k),
\]

то после сдвига всех индексов на \(h\) средний вектор не меняется, а
ковариационные элементы зависят от тех же разностей:

\[
(n_j+h)-(n_i+h)=n_j-n_i.
\]

Значит не меняется все конечномерное распределение.

\subsection{7. Что написать на
доске}\label{ux447ux442ux43e-ux43dux430ux43fux438ux441ux430ux442ux44c-ux43dux430-ux434ux43eux441ux43aux435}

\begin{enumerate}
\def\labelenumi{\arabic{enumi}.}
\tightlist
\item
  Нормальная величина:
\end{enumerate}

\[
X\sim N(m,\sigma^2),
\qquad
\varphi_X(t)=\exp\left\{itm-\frac12\sigma^2t^2\right\}.
\]

\begin{enumerate}
\def\labelenumi{\arabic{enumi}.}
\setcounter{enumi}{1}
\tightlist
\item
  Гауссовский вектор:
\end{enumerate}

\[
X\ \text{гауссовский}
\Longleftrightarrow
\langle u,X\rangle\ \text{нормальна для всех}\ u.
\]

\begin{enumerate}
\def\labelenumi{\arabic{enumi}.}
\setcounter{enumi}{2}
\tightlist
\item
  Характеристическая функция:
\end{enumerate}

\[
\varphi_X(u)
=
\exp\left\{i\langle u,m\rangle-\frac12u^T\Sigma u\right\}.
\]

\begin{enumerate}
\def\labelenumi{\arabic{enumi}.}
\setcounter{enumi}{3}
\tightlist
\item
  Ковариационная матрица:
\end{enumerate}

\[
\Sigma_{ij}=\operatorname{Cov}(X_i,X_j),
\qquad
u^T\Sigma u\ge0.
\]

\begin{enumerate}
\def\labelenumi{\arabic{enumi}.}
\setcounter{enumi}{4}
\tightlist
\item
  Некоррелированность и независимость:
\end{enumerate}

\[
\text{для совместно гауссовских: }
\operatorname{Cov}(X,Y)=0
\Longleftrightarrow
X,Y\ \text{независимы}.
\]

\begin{enumerate}
\def\labelenumi{\arabic{enumi}.}
\setcounter{enumi}{5}
\tightlist
\item
  Гауссовская последовательность:
\end{enumerate}

\[
(X_{n_1},\ldots,X_{n_k})\ \text{гауссовский для всех конечных наборов}.
\]

\begin{enumerate}
\def\labelenumi{\arabic{enumi}.}
\setcounter{enumi}{6}
\tightlist
\item
  Для Гаусса:
\end{enumerate}

\[
\text{широкая стационарность}
\Longrightarrow
\text{узкая стационарность}.
\]

\subsection{8. Типичные дополнительные
вопросы}\label{ux442ux438ux43fux438ux447ux43dux44bux435-ux434ux43eux43fux43eux43bux43dux438ux442ux435ux43bux44cux43dux44bux435-ux432ux43eux43fux440ux43eux441ux44b}

\textbf{Что является главным объектом билета?}

Нормальные случайные величины, гауссовские векторы и гауссовские
последовательности. Главное отличие гауссовского объекта: все
конечномерные распределения задаются средними и ковариациями.

\textbf{Почему в определении гауссовского вектора нужны все линейные
комбинации?}

Потому что нормальность отдельных координат не определяет совместный
закон. Нужно требовать, чтобы

\[
\sum_j u_jX_j
\]

была нормальной для любых коэффициентов \(u_j\).

\textbf{Что задает закон гауссовского вектора?}

Средний вектор \(m\) и ковариационная матрица \(\Sigma\).
Характеристическая функция:

\[
\varphi_X(u)
=
\exp\left\{i\langle u,m\rangle-\frac12u^T\Sigma u\right\}.
\]

\textbf{Какие условия нужны на \(\Sigma\)?}

Матрица должна быть симметричной и неотрицательно определенной:

\[
u^T\Sigma u\ge0.
\]

Если \(\Sigma\) положительно определена, есть плотность в
\(\mathbb R^n\); если вырождена, распределение сосредоточено на
подпространстве.

\textbf{Когда некоррелированность означает независимость?}

Для совместно гауссовских величин или групп координат. В общем случае
это неверно.

\textbf{Почему для гауссовских некоррелированность дает независимость?}

Потому что при нулевых межгрупповых ковариациях ковариационная матрица
становится блочно-диагональной, и характеристическая функция
раскладывается в произведение характеристических функций групп.

\textbf{Что такое гауссовская последовательность?}

Это последовательность, у которой любой конечный вектор

\[
(X_{n_1},\ldots,X_{n_k})
\]

гауссовский.

\textbf{Как задать гауссовскую последовательность?}

Нужно задать среднюю функцию \(m_n\) и неотрицательно определенную
ковариационную функцию \(K(n,m)\). Тогда конечномерные гауссовские
распределения согласованы и задают процесс.

\textbf{Почему широкая стационарность гауссовской последовательности
влечет узкую?}

Потому что гауссовские конечномерные распределения полностью
определяются средними и ковариациями. При широкой стационарности они не
меняются при сдвиге.

\textbf{Какая главная ошибка в этом билете?}

Считать, что если каждая координата нормальна, то весь вектор
автоматически гауссовский. Это неверно: нужна нормальность всех линейных
комбинаций.

\subsection{9. Типичные
ошибки}\label{ux442ux438ux43fux438ux447ux43dux44bux435-ux43eux448ux438ux431ux43aux438}

\begin{itemize}
\tightlist
\item
  Считать нормальность координат достаточной для совместной гауссовости.
\item
  Забывать вырожденные нормальные распределения.
\item
  Писать плотность многомерного нормального закона при вырожденной
  \(\Sigma\).
\item
  Забывать условие \(\Sigma\ge0\).
\item
  Переносить свойство ``некоррелированность равна независимости'' на
  негауссовские распределения.
\item
  Путать ковариационную матрицу и корреляционную матрицу.
\item
  Считать, что широкая стационарность всегда влечет узкую; это верно не
  вообще, а для гауссовских последовательностей.
\end{itemize}

\subsection{10. Связи с другими
билетами}\label{ux441ux432ux44fux437ux438-ux441-ux434ux440ux443ux433ux438ux43cux438-ux431ux438ux43bux435ux442ux430ux43cux438}

\begin{itemize}
\tightlist
\item
  Билет 08: характеристические функции задают распределения.
\item
  Билет 13: средние и ковариации получаются из производных
  характеристической функции.
\item
  Билет 14: ковариация, корреляция и некоррелированность; для Гаусса
  некоррелированность дает независимость.
\item
  Билет 15: согласованные гауссовские конечномерные распределения задают
  гауссовскую последовательность.
\item
  Билет 17: для гауссовских последовательностей широкая стационарность
  влечет узкую.
\item
  Билет 24: корреляционная функция стационарной гауссовской
  последовательности.
\item
  Билет 25: спектральные характеристики стационарных последовательностей
  второго порядка.
\end{itemize}

\subsection{11. Минимум на
удовлетворительно}\label{ux43cux438ux43dux438ux43cux443ux43c-ux43dux430-ux443ux434ux43eux432ux43bux435ux442ux432ux43eux440ux438ux442ux435ux43bux44cux43dux43e}

Нужно сказать:

\begin{enumerate}
\def\labelenumi{\arabic{enumi}.}
\tightlist
\item
  Одномерный нормальный закон:
\end{enumerate}

\[
X\sim N(m,\sigma^2),
\qquad
\varphi_X(t)=\exp\left\{itm-\frac12\sigma^2t^2\right\}.
\]

\begin{enumerate}
\def\labelenumi{\arabic{enumi}.}
\setcounter{enumi}{1}
\tightlist
\item
  Гауссовский вектор определяется нормальностью всех линейных
  комбинаций:
\end{enumerate}

\[
\langle u,X\rangle\ \text{нормальна для всех}\ u.
\]

\begin{enumerate}
\def\labelenumi{\arabic{enumi}.}
\setcounter{enumi}{2}
\tightlist
\item
  Его характеристическая функция:
\end{enumerate}

\[
\varphi_X(u)=
\exp\left\{i\langle u,m\rangle-\frac12u^T\Sigma u\right\}.
\]

\begin{enumerate}
\def\labelenumi{\arabic{enumi}.}
\setcounter{enumi}{3}
\item
  Гауссовская последовательность: каждый конечный вектор гауссовский.
\item
  Главная ловушка: нормальность координат не равна совместной
  гауссовости.
\end{enumerate}

\subsection{12. Минимум на
хорошо}\label{ux43cux438ux43dux438ux43cux443ux43c-ux43dux430-ux445ux43eux440ux43eux448ux43e}

Помимо минимума, нужно объяснить:

\begin{itemize}
\tightlist
\item
  что такое вырожденный нормальный закон;
\item
  почему \(\Sigma\) должна быть неотрицательно определенной;
\item
  почему линейное преобразование гауссовского вектора гауссово;
\item
  почему для совместно гауссовских некоррелированность равна
  независимости;
\item
  как задается гауссовская последовательность через \(m_n\) и
  \(K(n,m)\);
\item
  почему для гауссовских широкая стационарность влечет узкую.
\end{itemize}

\subsection{13. Что добавить для
отлично}\label{ux447ux442ux43e-ux434ux43eux431ux430ux432ux438ux442ux44c-ux434ux43bux44f-ux43eux442ux43bux438ux447ux43dux43e}

Для сильного ответа стоит добавить:

\begin{itemize}
\tightlist
\item
  плотность многомерного нормального закона при \(\Sigma>0\);
\item
  предупреждение про вырожденный случай \(\det\Sigma=0\);
\item
  доказательство через характеристические функции;
\item
  пример нормальных координат без совместной гауссовости;
\item
  блочно-диагональную ковариационную матрицу как критерий независимости
  групп;
\item
  связь с теоремой Колмогорова для построения гауссовской
  последовательности;
\item
  связь со стационарностью и спектральными характеристиками.
\end{itemize}

\subsection{14. Устная версия ответа на 5
минут}\label{ux443ux441ux442ux43dux430ux44f-ux432ux435ux440ux441ux438ux44f-ux43eux442ux432ux435ux442ux430-ux43dux430-5-ux43cux438ux43dux443ux442}

\begin{enumerate}
\def\labelenumi{\arabic{enumi}.}
\tightlist
\item
  Начинаю с одномерного нормального закона:
\end{enumerate}

\[
X\sim N(m,\sigma^2),
\qquad
\varphi_X(t)=\exp\left\{itm-\frac12\sigma^2t^2\right\}.
\]

\begin{enumerate}
\def\labelenumi{\arabic{enumi}.}
\setcounter{enumi}{1}
\tightlist
\item
  Гауссовский вектор \(X=(X_1,\ldots,X_n)\) определяется так: любая
  линейная комбинация координат нормальна.
\item
  Это сильнее, чем нормальность каждой координаты отдельно.
\item
  Средний вектор и ковариационная матрица:
\end{enumerate}

\[
m=\mathbb EX,
\qquad
\Sigma_{ij}=\operatorname{Cov}(X_i,X_j).
\]

\begin{enumerate}
\def\labelenumi{\arabic{enumi}.}
\setcounter{enumi}{4}
\tightlist
\item
  Характеристическая функция:
\end{enumerate}

\[
\varphi_X(u)
=
\exp\left\{i\langle u,m\rangle-\frac12u^T\Sigma u\right\}.
\]

\begin{enumerate}
\def\labelenumi{\arabic{enumi}.}
\setcounter{enumi}{5}
\tightlist
\item
  Поэтому гауссовский закон полностью задается \(m\) и \(\Sigma\).
\item
  Линейное преобразование гауссовского вектора снова гауссово:
\end{enumerate}

\[
Y=AX+b.
\]

\begin{enumerate}
\def\labelenumi{\arabic{enumi}.}
\setcounter{enumi}{7}
\tightlist
\item
  Для совместно гауссовских величин нулевая ковариация означает
  независимость. В общем случае это неверно.
\item
  Гауссовская последовательность - это последовательность, у которой
  любой конечный вектор гауссовский.
\item
  Она задается средней функцией \(m_n\) и ковариационной функцией
  \(K(n,m)\).
\item
  Для гауссовской последовательности широкая стационарность влечет
  узкую, потому что конечномерные распределения задаются средними и
  ковариациями.
\item
  Примеры: стандартная нормальная величина, белый гауссовский шум,
  гауссовский \(AR(1)\).
\end{enumerate}

\subsection{15. Если осталось 2 минуты перед
экзаменом}\label{ux435ux441ux43bux438-ux43eux441ux442ux430ux43bux43eux441ux44c-2-ux43cux438ux43dux443ux442ux44b-ux43fux435ux440ux435ux434-ux44dux43aux437ux430ux43cux435ux43dux43eux43c}

Запомнить главное.

Одномерная нормальная:

\[
\varphi_X(t)=\exp\left\{itm-\frac12\sigma^2t^2\right\}.
\]

Гауссовский вектор:

\[
\langle u,X\rangle\ \text{нормальна для всех}\ u.
\]

Характеристическая функция:

\[
\varphi_X(u)
=
\exp\left\{i\langle u,m\rangle-\frac12u^T\Sigma u\right\}.
\]

Гауссовская последовательность:

\[
(X_{n_1},\ldots,X_{n_k})\ \text{гауссовский для всех конечных наборов}.
\]

Главные ловушки:

\begin{itemize}
\tightlist
\item
  нормальность координат не означает совместную гауссовость;
\item
  некоррелированность равна независимости только в совместно гауссовском
  случае;
\item
  плотность с \(\Sigma^{-1}\) пишется только при невырожденной
  \(\Sigma\).
\end{itemize}

\newpage

\end{document}
